\documentclass[border=10pt]{standalone}
\usepackage{fontspec}
\AtBeginDocument{\hyphenation{zenodo}}% never hyphenate "zenodo" inside a DOI
\setmainfont{Noto Serif}[Ligatures=TeXOff]
\XeTeXgenerateactualtext=1
\newfontfamily\bengalifont{Noto Serif Bengali}[Script=Bengali, NFSSFamily=bnfam]
\usepackage[Latin,Bengali]{ucharclasses}
% group-free font switching (as in header-common.tex): robust inside braces, cells and quotes
\makeatletter
\def\BN@saved{\rmdefault}
\def\BN@bn{bnfam}
\newcommand\BNin{\edef\BN@tmp{\f@family}\ifx\BN@tmp\BN@bn\else\let\BN@saved\BN@tmp\fi\fontfamily{bnfam}\selectfont}
\newcommand\BNout{\fontfamily{\BN@saved}\selectfont}
\makeatother
\setTransitionsFor{Bengali}{\BNin}{\BNout}
% Typographic quotes and dashes are not in a ucharclasses class (class 0), so a Bangla word right
% after “ or ‘ would not trigger the Bangla font. Treat them as Latin so that every switch fires.
\XeTeXcharclass"2018=\XeTeXcharclass`A
\XeTeXcharclass"2019=\XeTeXcharclass`A
\XeTeXcharclass"201C=\XeTeXcharclass`A
\XeTeXcharclass"201D=\XeTeXcharclass`A
\XeTeXcharclass"2013=\XeTeXcharclass`A
\XeTeXcharclass"2014=\XeTeXcharclass`A
\XeTeXcharclass"2026=\XeTeXcharclass`A
\usepackage{newunicodechar}
\newunicodechar{।}{{\bengalifont\char"0964}}
\newunicodechar{॥}{{\bengalifont\char"0965}}

\usepackage[table]{xcolor}
\usepackage{booktabs,tabularx,array}
\definecolor{bapilrule}{HTML}{1F2937}
\definecolor{bapilstripe}{HTML}{F3F4F6}
\renewcommand{\arraystretch}{1.28}
\setlength{\tabcolsep}{5pt}

\begin{document}
\begin{minipage}{17cm}
{\large\bfseries সারণি ৬৫।}\ {\large ৩৬ মাসের পরিকল্পনার আগে ও শুরুতে প্রয়োজনীয় পূর্বশর্ত ও অনুমোদন}\par\vspace{7pt}
\small
\rowcolors{2}{bapilstripe}{white}
\begin{tabularx}{\linewidth}{@{}>{\hsize=1.170\hsize\raggedright\arraybackslash}X>{\hsize=1.011\hsize\raggedright\arraybackslash}X>{\hsize=1.383\hsize\raggedright\arraybackslash}X>{\hsize=0.798\hsize\raggedright\arraybackslash}X>{\hsize=0.638\hsize\raggedright\arraybackslash}X@{}}
\arrayrulecolor{bapilrule}\toprule
\rowcolor{white}\textbf{পূর্বশর্ত} & \textbf{কেন প্রয়োজন} & \textbf{দলিল বা কর্তৃপক্ষ} & \textbf{অবস্থা (অক্টোবর ২০২৬)} & \textbf{কখন প্রয়োজন} \\
\midrule
ধাপ ০-এর অর্থায়ন ও নির্দেশ & যেকোনো প্রতিশ্রুতির আগে স্বাগতিকের আহ্বান, আইনি পর্যালোচনা, ব্যয়-নিরূপণ সমীক্ষা ও ডেটা প্রবেশাধিকার নিশ্চিতকরণ & সিদ্ধান্ত D5 (আইসিটি বিভাগ বা কোনো উন্নয়ন সহযোগী); প্রায় ০.৪৪ কোটি টাকা (অনুচ্ছেদ ১৭.২) & অনুরোধ করা হয়নি & মাস ১-এর প্রায় ছয় মাস আগে \\
স্বাগতিক প্রতিষ্ঠান নির্বাচন & জনবল, অর্থ ও ডেটার আইনগত আবাস & ধাপ ০-এ উন্মুক্ত আহ্বান (অনুচ্ছেদ ৫.২); স্বাগতিকের পরিচালনা পর্ষদের সিদ্ধান্ত & শুরু হয়নি & ধাপ ০-এর শেষে \\
আইসিটি বিভাগের সঙ্গে সমঝোতা স্মারক & নীতি-যোগাযোগ, ডেটা প্রবেশাধিকার ও হস্তান্তরের পথ & আইসিটি বিভাগের নিজস্ব অনুমোদন-প্রক্রিয়ায় স্বাক্ষরিত সমঝোতা স্মারক (সিদ্ধান্ত D2) & এই ব্লুপ্রিন্টে প্রস্তাবিত; সিদ্ধান্ত হয়নি & মাস ১ (সময়গণনা শুরু) \\
প্রথম বছরের অর্থায়ন নিশ্চিত & বেতন, যন্ত্রপাতি ও সমীক্ষা & প্রযোজ্য অনুমোদন-প্রক্রিয়ায় সরকারি বাজেট বরাদ্দ এবং অনধিক ৩০\%-এর জন্য উন্নয়ন সহযোগীর চুক্তি (বিদেশি অনুদান অর্থনৈতিক সম্পর্ক বিভাগের মাধ্যমে); গবেষণা অনুদান ও শিল্পখাতের অর্থ কেবল দ্বিতীয় বছর থেকে (অনুচ্ছেদ ১৭.৪) & কোনো অর্থ চাওয়া বা অনুমোদিত হয়নি & মাস ১ \\
স্বাগতিক-চুক্তি ও পৃথক ব্যয়কেন্দ্র & আর্থিক শাসন (অনুচ্ছেদ ১৭.৮) & স্বাগতিকের আর্থিক বিধি & শুরু হয়নি & মাস ১ \\
পরিচালক নিয়োগ & জবাবদিহির একক কেন্দ্র & স্বাগতিকের নিয়োগবিধি; স্টিয়ারিং বোর্ডের নিশ্চিতকরণ & শুরু হয়নি & মাস ৩ \\
স্টিয়ারিং বোর্ড ও ERRC গঠন; সনদ গ্রহণ & সিদ্ধান্তের অধিকার ও সুরক্ষাব্যবস্থা & মনোনয়ন (সিদ্ধান্ত D4); সনদ (সংযুক্তি ঙ) & শুধু খসড়া সনদ & মাস ৩ (MS1) \\
উন্মুক্ত নথির পুনর্ব্যবহারের শর্ত ও একসঙ্গে প্রবেশাধিকার & P1 কর্পাস & ডেটা-রক্ষকদের লিখিত নিশ্চয়তা (সিদ্ধান্ত D3); বিকল্প হিসেবে তথ্য অধিকার আবেদন & শুরু হয়নি & ধাপ ০-এ অনুরোধ; মাস ৬-এর মধ্যে নিশ্চিত \\
স্বাধীন আইনি পর্যালোচনা & কোন দায়িত্ব কবে থেকে প্রযোজ্য, এবং স্বাগতিক ERRC-র ভেটো মানতে পারে কি না, তা নিশ্চিত করা & ধাপ ০-এ নিযুক্ত আইনজীবী (অনুচ্ছেদ ১৪.৩) & শুরু হয়নি & ধাপ ০-এ, স্বাগতিক-চুক্তি স্বাক্ষরের আগে \\
বাংলাদেশে হোস্টিং & প্রয়োজনীয়তা NFR-02 & জাতীয় ডেটা সেন্টার বা স্বাগতিকের ডেটা সেন্টারের সঙ্গে চুক্তি & শুরু হয়নি & মাস ৪ \\
যন্ত্রপাতি ক্রয় & পাইলট অবকাঠামো & ধাপ ০-এর ব্যয়-নিরূপণ সমীক্ষার দরপ্রস্তাব ব্যবহার করে স্বাগতিকের সরকারি ক্রয়বিধি; শুল্ক ও ভ্যাট অব্যাহতি থাকলে তার নিশ্চিতকরণ & শুরু হয়নি & মাস ৩ থেকে ৬ \\
গবেষণা-নৈতিকতার অনুমোদন & ব্যবহারকারী-সমীক্ষা, P3 ও P6 & স্বাগতিকের গবেষণা-নৈতিকতা কমিটি; ERRC & শুরু হয়নি & প্রতিটি সমীক্ষার আগে \\
অ-উন্মুক্ত ডেটার জন্য ডেটা-বিনিময় চুক্তি & P3 & চুক্তি (সংযুক্তি ঝ); DPIA (সংযুক্তি জ); সম্ভব হলে NRDEX & শুরু হয়নি & মাস ১২ থেকে \\
এআই হাব প্রকল্পের সঙ্গে কর্মবিভাজন-নোট & দ্বৈততা এড়ানো & আইসিটি বিভাগ ও বাংলাদেশ হাই-টেক পার্ক কর্তৃপক্ষের সঙ্গে পত্র বা সমঝোতা স্মারক & শুরু হয়নি & MS2 \\
\bottomrule
\end{tabularx}
\par\vspace{5pt}{\footnotesize\itshape সিদ্ধান্ত D5-এর পর ধাপ ০ প্রায় ছয় মাস চলে। পরিকল্পনার মাস ১ হলো প্রথম সেই মাস, যখন সমঝোতা স্মারক, স্বাগতিক-চুক্তি ও প্রথম বছরের অর্থায়ন তিনটিই নিশ্চিত। ২০২৬-এর চতুর্থ ত্রৈমাসিকে D5 গৃহীত হলে ধাপ ০ ২০২৭-এর প্রথমার্ধে চলতে পারে এবং মাস ১ হতে পারে জুলাই ২০২৭, অর্থাৎ ২০২৭-২৮ অর্থবছরের শুরু, যাতে প্রথম বছরের অর্থ সেই বছরের বাজেটে রাখা যায়।}
\par\vspace{3pt}{\scriptsize\color{gray} উৎস: BAPIL টেকনিক্যাল ব্লুপ্রিন্ট, সংস্করণ ১.০ (২০২৬)। CC BY 4.0।}
\end{minipage}
\end{document}
